\chapter*{Préambule}
\label{preambule}
\addcontentsline{toc}{chapter}{Préambule}
\markboth{Préambule}{Préambule}

J'ai écrit ce livre pour aider les étudiants en fin de cycle à se familiariser avec un outil qui leur permette d'analyser leurs données.

Lorsque j'étais étudiant en médecine, arrivé en fin de formation, j'ai été confronté à un problème épineux : l'analyse des données recueillies pour ma thèse de doctorat en médecine. Dans notre unité de formation, Epi Info était de très loin le logiciel le plus utilisé à cet effet. J'ai donc entrepris de me familiariser avec cet outil.

Ce livre est né d'un constat : les ressources documentaires en français étaient très limitées, même s'il existe un tutoriel rédigé par les concepteurs du logiciel.

\section*{Pourquoi le logiciel Epi Info ?}

Epi Info est un logiciel gratuit créé par les \emph{Centers for Disease Control and Prevention} (CDC), en français Centres pour le contrôle et la prévention des maladies, aux États-Unis. Il permet de concevoir des formulaires de saisie de données et de réaliser des analyses statistiques. Il est téléchargeable gratuitement sur le site du CDC. Sa courbe d'apprentissage est courte : même un néophyte peut le prendre en main très facilement.

Il faut toutefois noter quelques inconvénients. Epi Info 7 n'est pas multiplateforme : il a été conçu pour le système d'exploitation Windows de Microsoft\footnote{Notons néanmoins que Windows est le système d'exploitation le plus répandu sur les ordinateurs de bureau.}. Pour les représentations graphiques, bien que le logiciel le permette, nous les réaliserons en général dans un logiciel tiers comme Microsoft Excel ou LibreOffice Calc. De même, pour des analyses plus poussées, vous aurez besoin de logiciels plus puissants.

\section*{Organisation du livre}

Le livre suit les étapes d'une enquête, de la conception du questionnaire à l'analyse des données :
\begin{itemize}
	\item le chapitre~\ref{installationConfig} présente l'installation d'Epi Info 7 et son passage en français ;
	\item le chapitre~\ref{creerQuestionnaire} montre comment créer un questionnaire avec ses différents types de champs ;
	\item le chapitre~\ref{codeValidation} introduit le code de contrôle, qui rend le questionnaire plus robuste ;
	\item le chapitre~\ref{saisirDonnees} traite de la saisie des données ;
	\item les chapitres~\ref{moduleTableauDeBord} et~\ref{moduleClassic} sont consacrés à l'analyse statistique avec les modules \module{tableau de bord visuel} et \module{classique} ;
	\item le chapitre~\ref{moduleStatcalc} présente le module \module{statcalc}, qui permet de calculer la taille d'un échantillon et d'analyser rapidement des tableaux déjà constitués.
\end{itemize}

Au fil de ce livre, vous apprendrez progressivement à créer un formulaire robuste et cohérent, puis à analyser vos données avec ce logiciel.

Je vous souhaite une bonne lecture.
